\documentclass[
	aspectratio=169, % 16:9 slides. Use 'aspectratio=43' for classic 4:3.
]{beamer}

\usetheme{scribe} % The scribe beamer theme (matches scribe.cls)
% Theme options:
%   \usetheme[noframenumbers]{scribe} % hide the frame counter
%   \usetheme[nosectionpages]{scribe} % no automatic section divider slides

\usepackage{lipsum} % For generating dummy text

% Define different aliases (along with colors) for comments
\scribedefinecomment{fb}{orange}
\scribedefinecomment{jd}{blue}

% Uncomment the following line to hide comments in the final slides.
% \scribeshowcommentsfalse

\title{An example presentation\\ using the scribe beamer theme}
\subtitle{Minimalist slides, matching the paper}

\author{Federico Bruzzone \orcidlink{0000-0002-8701-8853}}
\institute{%
	Computer Science Department, Università degli Studi di Milano\\
	\href{mailto:federico.bruzzone@unimi.it}{\texttt{federico.bruzzone@unimi.it}}
}
\date{May 2026}

% Short versions used in the footline
\title[scribe beamer]{An example presentation\\ using the scribe beamer theme}
\author[F. Bruzzone]{Federico Bruzzone \orcidlink{0000-0002-8701-8853}}

\begin{document}

% ============================
% ========== Title ===========
% ============================
\begin{frame}[plain,noframenumbering]
	\titlepage
\end{frame}

% ============================
% ========== Outline =========
% ============================
\begin{frame}{Outline}
	\tableofcontents
\end{frame}

% ============================
% ========== Content =========
% ============================
\section{Introduction}

\begin{frame}[fragile]{A clean first slide}
	The \tool{scribe} beamer theme reuses the paper's identity: Libertine
	serif, Bera Mono, and the \textcolor{scribeaccent}{MidnightBlue} accent.

	\begin{itemize}
		\item Minimalist, content-first layout~\cite{Scribe25}
		\item A thin accent rule under every frame title
		\item Triangle bullets in the accent color
		      \begin{itemize}
			      \item Second-level items inherit the style
		      \end{itemize}
	\end{itemize}

	\fbcomment{This is a comment in orange from Federico Bruzzone.}

	\scribedivider

	The \verb|\scribedivider| ornament works on slides too.
\end{frame}

\begin{frame}{Blocks and dingbats}
	\begin{columns}[T]
		\begin{column}{0.5\textwidth}
			\begin{block}{Definition}
				A block with the accent-tinted body.
			\end{block}
			\begin{alertblock}{Warning}
				An alert block for caveats.
			\end{alertblock}
			\begin{exampleblock}{Example}
				An example block for, well, examples.
			\end{exampleblock}
		\end{column}
		\begin{column}{0.5\textwidth}
			Colored dingbats, ported from \tool{scribe}:
			\begin{itemize}
				\item \dinggreen{52}~done
				\item \dingred{56}~failed
				\item \dingblue{43}~note
			\end{itemize}
			\jdcomment{Another comment, in blue, from John Doe.}
		\end{column}
	\end{columns}
\end{frame}

\section{Math and code}

\begin{frame}{Mathematics}
	The factorial, set with the Libertine math defaults:
	\begin{equation*}
		fact(n) = \begin{cases}
			1                 & \text{if } n = 0 \\
			n \cdot fact(n-1) & \text{if } n > 0
		\end{cases}
	\end{equation*}
\end{frame}

\begin{frame}[fragile]{Code listings}
	Monospace is set in Bera Mono:
\begin{lstlisting}[basicstyle=\ttfamily\small, keywordstyle=\color{scribeaccent}\bfseries, language=C]
int fact(int n) {
    if (n == 0) return 1;
    return n * fact(n - 1);
}
\end{lstlisting}
\end{frame}

\section{Wrap-up}

\begin{frame}[fragile]{Summary}
	\begin{itemize}
		\item One theme file: \texttt{beamerthemescribe.sty}
		\item Same fonts, colors, and ornament as \texttt{scribe.cls}
		\item Toggle comments with \verb|\scribeshowcommentsfalse|
	\end{itemize}
\end{frame}

\begin{frame}[plain,noframenumbering]
	\centering
	\vfill
	{\Large\bfseries Thank you!}
	\scribedivider
	\href{https://federicobruzzone.github.io/}{\texttt{federicobruzzone.github.io}}
	\vfill
\end{frame}

% ============================
% ========== Refs ============
% ============================
\begin{frame}[allowframebreaks]{References}
	\bibliographystyle{unsrt}
	\bibliography{local}
\end{frame}

\end{document}
